\BourbakiExerciseGroup

\begin{enumerate}
\def\labelenumi{\arabic{enumi})}
\tightlist
\item
  A point \(a \in \mathbb{R}\) is said to be left adherent to a subset \(\mathbf{A}\) of \(\mathbb{R}\) if it lies in the closure of \(\mathbf{A}\) and if there is an interval \(]a, b[ (a < b)\) not containing any point of \(\mathbf{A}\) . Show that the set of left adherent points of a subset \(\mathbf{A}\) of \(\mathbb{R}\) is countable (set up a one-to-one correspondence between the set of left adherent points and a set of pairwise disjoint open intervals). Hence show that every well-ordered subset of \(\mathbb{R}\) is countable.
\end{enumerate}

¶ 2) A countable dense subset \(\mathbf{A}\) of \(\mathbb{R}\) is not closed. {[}If \((a_{n})\) is a sequence obtained by arranging the points of \(\mathbf{A}\) in some order, define a sequence of intervals \([b_{n}, c_{n}]\) such that \(b_{n-1} < b_{n} < c_{n} < c_{n-1}\) for all \(n\) , and such that \([b_{n}, c_{n}]\) contains no point \(a_{k}\) with index \(k \leqslant n\) ; then use Theorem 2.{]} Deduce that \(\mathbb{R}\) is uncountable (cf.~§ 8, no. 6).

\begin{enumerate}
\def\labelenumi{\arabic{enumi})}
\setcounter{enumi}{2}
\item
  Let \((\mathbf{I}_n)\) be an infinite sequence of non-empty open intervals in \(\mathbf{R}\) such that \(\overline{\mathbf{I}}_n \cap \overline{\mathbf{I}}_m = \emptyset\) for \(m \neq n\) . Show that the complement of the union of the \(\mathbf{I}_n\) is a perfect set (Chapter I, § 2, no. 6); in particular, Cantor's triadic set is perfect.
\item
  The sum of the lengths of the intervals contiguous to Cantor's triadic set is equal to 1. Define a totally disconnected perfect subset of {[}0, 1{]} such that the sum of the lengths of the intervals contiguous to A and contained in {[}0, 1{]} is a given number m such that \(0 < m \leqslant 1\) .
\item
  Let \(l\) be a number \(> 0\) , and suppose that to each point \(x \in \mathbf{R}\) corresponds an open interval \(\mathbf{I}(x)\) with centre \(x\) and length \(\leqslant l\) . Show that every compact interval \([a, b]\) can be covered by a finite number of intervals \(\mathbf{I}(x_i)\) , the sum of whose lengths is \(\leqslant l + 2(b - a)\) . (Prove that if the proposition is true for each interval \([a, x]\) such that \(a \leqslant x < c\) , then there exists \(d > c\) such that it is true for each interval \([a, y]\) such that \(a \leqslant y < d\) .) If \(l = (b - a)/n\) , where \(n\) is an integer \(\geqslant 1\) , show that the result cannot be improved.
\end{enumerate}

¶ 6) a) Let X be a non-empty linearly ordered set, endowed with the topology \(\mathcal{C}_{0}(X)\) (Chapter I, § 2, Exercise 5). Show that X is compact if and only if each subset of X has a least upper bound, in other words if and only if X is a complete lattice (Set Theory, Chapter III, § 1, Exercise 11). (To show that the condition is necessary, argue as for Theorem 3; for sufficiency, consider a filter \(\mathfrak{F}\) on X and show that, if A is the set of greatest lower bounds of sets of \(\mathfrak{F}\) , then the least upper bound of A is a cluster point of \(\mathfrak{F}\) ).

\begin{enumerate}
\def\labelenumi{\alph{enumi})}
\setcounter{enumi}{1}
\tightlist
\item
  Give an example of a complete lattice X which is not compact in the topology \(\mathcal{C}_{0}(\mathbf{X})\) .
\end{enumerate}

¶ 7) Let X be a non-empty linearly ordered set, endowed with the topology \(\mathfrak{G}_{\mathfrak{0}}(\mathbf{X})\) .

\begin{enumerate}
\def\labelenumi{\alph{enumi})}
\tightlist
\item
  If X is connected, show that it has the following two properties:
\end{enumerate}

α) Every non-empty subset of X which is bounded above has a least upper bound (if A is non-empty and bounded above, and if B is the set of upper bounds of A, and C the set of lower bounds of B, show that \(B \cup C = X\) and that B and C are closed).

β) The set X is without gaps, in other words (Set Theory, Chapter III, § 1, Exercise 18) every open interval {]}a, b{[} (a \textless{} b) of X is non-empty (use the argument of Theorem 4).

\begin{enumerate}
\def\labelenumi{\alph{enumi})}
\setcounter{enumi}{1}
\item
  Show conversely that if X has properties \(\alpha\) and \(\beta\) ), then X is connected. (Show first that every closed interval \([a, b]\) is connected: assuming that this interval is the disjoint union of two non-empty closed sets A and B, and that \(a \in A\) , consider the greatest lower bound of B and hence obtain a contradiction).
\item
  Show that if X is connected it is locally compact and locally connected, and that the only connected subsets of X are the intervals (bounded or unbounded).
\end{enumerate}

\begin{enumerate}
\def\labelenumi{\arabic{enumi})}
\setcounter{enumi}{7}
\tightlist
\item
  Let X and Y be two linearly ordered sets, endowed with the topologies \(\mathcal{G}_{0}(X)\) and \(\mathcal{G}_{0}(Y)\) respectively.
\end{enumerate}

\begin{enumerate}
\def\labelenumi{\alph{enumi})}
\tightlist
\item
  Suppose \(\mathbf{X}\) connected and let \(f: \mathbf{X} \to \mathbf{Y}\) be a continuous mapping. Show that, for all \(x, y\) in \(\mathbf{X}\) such that \(x < y\) , every \(z \in \mathbf{Y}\) belonging to the closed interval with end-points \(f(x)\) and \(f(y)\) belongs to \(f(\mathbf{X})\)
\end{enumerate}

(use Exercise 7). Deduce that \(f\) is a homeomorphism of \(\mathbf{X}\) onto \(f(\mathbf{X})\) if and only if \(f\) is continuous and strictly monotone.

\begin{enumerate}
\def\labelenumi{\alph{enumi})}
\setcounter{enumi}{1}
\tightlist
\item
  Give an example of a homeomorphism of the rational line Q onto itself which is not strictly monotone.
\end{enumerate}

\begin{enumerate}
\def\labelenumi{\arabic{enumi})}
\setcounter{enumi}{8}
\tightlist
\item
  \begin{enumerate}
  \def\labelenumii{\alph{enumii})}
  \tightlist
  \item
    Let X be a countable linearly ordered set without gaps (Exercise 7). Show that there is a strictly increasing bijective mapping \(\varphi\) of X onto one of the four intervals of Q with end-points o and i. {[}Arrange the elements of X and the points of the appropriate interval of Q in sequences \((a_n), (b_n)\) , and define \(\varphi\) by induction.{]} If X is endowed with the topology \(\tilde{C}_0(X)\) , then \(\varphi\) is a homeomorphism of X onto \(\varphi(X)\) .
  \end{enumerate}
\end{enumerate}

\begin{enumerate}
\def\labelenumi{\alph{enumi})}
\setcounter{enumi}{1}
\item
  Deduce from a) that every countable dense subset of an open interval {]}a, b{[} of R is homeomorphic to Q.
\item
  Deduce from a) that if X is any countable ordered set endowed with the topology \(\mathfrak{G}_{0}(\mathbf{X})\) , then there exists a strictly increasing homeomorphism of X onto a subspace of Q {[}embed X in a countable linearly ordered space \(X'\) without gaps, such that \(\mathfrak{G}_{0}(\mathbf{X})\) is induced by \(\mathfrak{G}_{0}(\mathbf{X}')\) {]}.
\end{enumerate}

\begin{enumerate}
\def\labelenumi{\arabic{enumi})}
\setcounter{enumi}{9}
\tightlist
\item
  Show that every countable subset of Cantor's triadic set K, which is dense in K and contains no end-point of an interval contiguous to K, is homeomorphic to the rational line Q (use Exercise 9).
\end{enumerate}

¶ 11) a) Let X be a linearly ordered set with the topology \(\mathfrak{G}_{0}(\mathbf{X})\) . If X is connected and has a countable dense subset A, show that there exists a homeomorphism (strictly monotonic by virtue of Exercise 8) of X onto one of the intervals of R with end-points o, i, which maps A onto the intersection of Q with this interval (use Exercises 9 and 7).

\begin{enumerate}
\def\labelenumi{\alph{enumi})}
\setcounter{enumi}{1}
\item
  Deduce from a) that if B is a subset of R whose complement is dense in R, then there exists a homeomorphism of R onto itself which maps B onto a subset of the set CQ of irrational numbers (consider a dense countable subset A of R contained in CB).
\item
  In the set \(\mathbf{R} \times \{0, 1\}\) , linearly ordered by the lexicographic ordering, let \(\mathbf{R}'\) denote the complement of \(\mathbf{Q} \times \{1\}\) . Show that \(\mathbf{R}'\) , endowed with the topology \(\mathfrak{G}_0(\mathbf{R}')\) , is locally compact and totally disconnected, and contains a countable dense subset \(\mathbf{Q}'\) homeomorphic to \(\mathbf{Q}\) , but that the topology induced on a non-empty open interval of \(\mathbf{R}'\) does not have a countable base; in particular such an interval cannot be homeomorphic to any subset of \(\mathbf{R}\) .
\end{enumerate}

¶ 12) a) Let A be a well-ordered set and let I denote the interval {[}0, 1{[} of R. Show that the set X = A × I, linearly ordered by the lexicographic ordering and endowed with the topology \(\mathfrak{G}_{0}(X)\) , is connected (Exercise 7); there are therefore linearly ordered sets X which are connected in the topology \(\mathfrak{G}_{0}(\mathbf{X})\) and have an arbitrary cardinal {[}cf.~§ 4, Exercise 7 b){]}.

\begin{enumerate}
\def\labelenumi{\alph{enumi})}
\setcounter{enumi}{1}
\tightlist
\item
  Take A to be an uncountable well-ordered set in which every segment \(\leftarrow, t]\) is countable: the corresponding topological space X is called the Alexandroff half-line. Show that for every interval \([a, b]\) of X (a \textless{} b) there exists a strictly increasing homeomorphism of \([a, b]\) onto the interval \([o, i]\) of R. (Prove by transfinite induction that there exists a strictly increasing homeomorphism of \(A \cap [a, b]\) , A being identified with the set of points \((t, o)\) of X onto a subspace of \([o, i]\) .)
\end{enumerate}

¶ 13) Let \(a\) be an irrational number \(>0\) . For each rational number \(x\) , let \(f_{a}(x)\) be the real number in the interval \([0, a]\) such that \(x - f_{a}(x)\) is an integral multiple of \(a\) . Show that \(f_{a}\) is a continuous injective mapping of \(\mathbf{Q}\) into \([0, a]\) . Deduce with the help of Exercise 9 that there exist continuous bijective mappings of \(\mathbf{Q}\) onto itself, whose inverse mapping is not continuous at any point.

¶ 14) Let X be a connected Hausdorff topological space not consisting of a single point, and let \(\Delta\) denote the diagonal of \(X \times X\) .

\begin{enumerate}
\def\labelenumi{\alph{enumi})}
\item
  Show that there is no partition of \(\mathbb{C}\Delta\) into two open sets \(\mathbf{D}_1, \mathbf{D}_2\) such that \(\mathbf{D}_1^{-1} = \mathbf{D}_1\) and \(\mathbf{D}_2^{-1} = \mathbf{D}_2\) . {[}Begin by remarking that \(\overline{\mathbf{D}_1(x)}\) and \(\overline{\mathbf{D}_2(x)}\) are connected for each \(x \in \mathbf{X}\) (use Exercise 4 a) of Chapter I, § 11); deduce that, if \(y \in \mathbf{D}_1(x)\) , we have \(\overline{\mathbf{D}_2(x)} \times \{y\} \subset \mathbf{D}_1\) , and show that this implies that, for each \(x \in \mathbf{X}\) , one or the other of the sets \(\mathbf{D}_1(x)\) , \(\mathbf{D}_2(x)\) is empty; hence show that one or the other of the sets \(\mathbf{D}_1, \mathbf{D}_2\) is empty, which is contrary to hypothesis.{]} Deduce that either \(\mathbb{C}\Delta\) is connected or else \(\mathbb{C}\Delta\) has exactly two components A, B such that \(\mathbf{B} = \overline{\mathbf{A}}\) .
\item
  A linear ordering on X is said to be compatible with the topology \(\mathfrak{C}\) of X if \(\mathfrak{C}_0(\mathbf{X})\) is coarser than \(\mathfrak{C}\) . Deduce from \(a\) ) that there exists such a linear ordering on X if and only if \(\mathbb{C}\Delta\) is not connected, and that there are then exactly two such linear orderings. {[}With the notation of \(a\) ), show that the order relation must be either \((x, y) \in \mathbf{A}\) or \((x, y) \in \mathbf{B} = \overline{\mathbf{A}}\) .{]} Show that the intervals, for these orderings, are connected in the topology \(\mathfrak{C}\) , and that they are the only connected subsets of X with respect to \(\mathfrak{C}\) .
\item
  Show that if X is either locally connected or locally compact in the topology \(\mathfrak{C}\) , and if there exists a linear ordering on X compatible with \(\mathfrak{C}\) , then we must have \(\mathfrak{C} = \mathfrak{C}_0(X)\) . {[}If X is locally connected, use b). If X is locally compact, show that every compact neighbourhood of a point \(x \in X\) is also a neighbourhood of \(x\) with respect to \(\mathfrak{C}_0(X)\) , by considering the frontier of the neighbourhood{]}.
\end{enumerate}

* d) Let X be the subspace of \(\mathbf{R}^2\) consisting of (o, o) and all pairs \((x, y)\) such that \(x > 0\) and \(y = \sin (\mathrm{i} / x)\) . Show that there exists a linear ordering on X which is compatible with the topology induced on X by the topology of \(\mathbf{R}^2\) , but that this topology is strictly finer than \(\mathcal{C}_0(X)\) . *

Give an example of a connected Hausdorff space X, no point of which has a fundamental system of connected neighbourhoods, and on which there is a linear ordering compatible with the topology of X {[}cf.~Chapter I, § 11, Exercise 2 b){]}.

¶ 15) a) Let X be a connected Hausdorff space such that, for each \(x \in X\) , \(\complement\{x\}\) has exactly two components. Show that, if K is either of these components, \(K = K \cup \{x\}\) {[}Chapter I, § 11, Exercise 4 a){]}. Let \(x, y\) be two distincts points of X, let A and B be the components of \(\complement\{x\}\) and let \(A', B'\) be the components of \(\complement\{y\}\) . Show that one of the two sets A, B is contained in \(A'\) or in \(B'\) .

\begin{enumerate}
\def\labelenumi{\alph{enumi})}
\setcounter{enumi}{1}
\tightlist
\item
  Let \(x_0\) be a point of \(X\) and let \(A(x_0), B(x_0)\) be the components of \(\complement\{x_0\}\) . For each \(x \neq x_0\) there is exactly one of the two components of \(\complement\{x\}\) which is either contained in \(A(x_0)\) or contains \(A(x_0)\) ; denote this component by \(A(x)\) . Show that the relation \(A(x) \subset A(y)\) is a linear order relation compatible (Exercise 14) with the topology of \(X\) . c) Extend the conclusion of b) to the case where \(\complement\{x\}\) has two components for all \(x \in X\) except for one point \(a \in X\) , for which \(\complement\{a\}\) is connected.
\end{enumerate}

* \(d\) ) Let \(\mathbf{X}\) be the subspace of \(\mathbf{R}^2\) consisting of the points \(a = (0, 1)\) , \(b = (0, -1)\) and all pairs \((x, y)\) such that \(x \neq 0\) and \(y = \sin(1/x)\) . Show that \(\mathbf{X}\) is connected in the topology \(\mathfrak{G}\) induced on \(\mathbf{X}\) by the topology of \(\mathbf{R}^2\) , that \(\mathbb{G}\{a\}\) and \(\mathbb{G}\{b\}\) are connected and that, for each point \(x\) of \(\mathbf{X}\) other than \(a\) and \(b\) , \(\mathbb{G}\{x\}\) has exactly two components; but that there exists no linear ordering on \(\mathbf{X}\) compatible with the topology \(\mathfrak{G}_{*}\) .

¶ 16) Let X be a connected Hausdorff space which is completely irreducible between two distinct points \(a\) and \(b\) , that is to say such that there is no connected subset of X other than X itself which contains both \(a\) and \(b\) . Show that \(\mathbf{X}' = \mathbb{C}\{a, b\}\) is connected and that, for each \(x \in \mathbf{X}'\) , the complement of \(x\) in X has exactly two components \(\mathbf{A}(x)\) and \(\mathbf{B}(x)\) such that \(a \in \mathbf{A}(x)\) , \(b \in \mathbf{B}(x)\) , \(a \notin \mathbf{B}(x)\) , \(b \notin \mathbf{A}(x)\) . {[}Show that \(\mathbb{C}\{x\}\) cannot be partitioned into three non-empty open sets, by using Exercise 4 a) of Chapter I, § 11.{]} Deduce that there exists a linear ordering on X which is compatible with the topology of X (use Exercise 15).

¶ 17) Let X be a connected Hausdorff space such that whenever X is covered by three non-empty connected sets distinct from X, there are two of these sets whose union is distinct from X.

\begin{enumerate}
\def\labelenumi{\alph{enumi})}
\item
  Show that, for every point \(x \in \mathbf{X}\) , \(\mathbb{G}\{x\}\) has at most two components (same method as in Exercise 16).
\item
  Show that there are at most two points of X whose complements are connected. Moreover, if there are two distinct points a, b with this property, show that, for each x other than a or b, a and b belong to different components of \(C\{x\}\) .
\item
  Deduce from a) and b) that there is a linear ordering on X compatible with the topology of X (use Exercise 15).
\end{enumerate}

¶ 18) a) Let X be a compact, connected, locally connected space which is irreducible between two of its points \(a\) , \(b\) (Chapter II, § 4, Exercise 19). Let \(x\) be a point of X other than \(a\) or \(b\) . Show that \(\mathbb{C}\{x\}\) has exactly two components and that \(a\) and \(b\) belong to different components of \(\mathbb{C}\{x\}\) . {[}Arguing as in Exercise 16, show first that \(\mathbb{C}\{x\}\) cannot have more than two components. Next, if V is any connected closed neighbourhood of \(x\) which contains neither \(a\) nor \(b\) , let \(\mathrm{A}_{\mathrm{V}}\) , \(\mathrm{B}_{\mathrm{V}}\) be the components of \(a\) , \(b\) respectively in \(\mathbb{C}\mathrm{V}\) ; show that \(\mathrm{A}_{\mathrm{V}} \neq \mathrm{B}_{\mathrm{V}}\) by using Exercise 17 of Chapter II, § 4; finally consider the union of the sets \(\mathrm{A}_{\mathrm{V}}\) (resp. \(\mathrm{B}_{\mathrm{V}}\) ) as V runs through the set of all connected closed neighbourhoods of \(x\) in X.{]} Deduce that there exists a linear ordering on X such that \(\mathfrak{C}_{0}(\mathbf{X})\) is the given topology on X {[}Exercises 15 and 14 c){]}.

\begin{enumerate}
\def\labelenumi{\alph{enumi})}
\setcounter{enumi}{1}
\tightlist
\item
  Let X be a compact connected space which is irreducible between two of its points a, b. Show that if X has more than one prime constituent (Chapter II, § 4, Exercise 20) then the space \(X'\) of prime constituents of X (loc. cit.) is irreducible between the prime constituents of a and b, and that there is a linear ordering on \(X'\) such that the topology of \(X'\) is identical with \(\mathcal{G}_{0}(\mathbf{X}')\) . * Give an example in which \(X'\) is distinct from X {[}cf.~Exercise 15 d{]}. *
\end{enumerate}

¶ 19) a) Let X be a connected Hausdorff space and let \(a, b\) be two distinct points of X such that there is no connected closed subset of X, other than X itself, which contains both \(a\) and \(b\) . If \(x\) is a point of X other than \(a\) or \(b\) , then \(\mathbb{C}\{x\}\) is not connected if and only if, for each \(y \in \mathbb{C}\{x\}\) , there is in \(\mathbb{C}\{y\}\) a connected subset which is closed in X and contains \(x\) and one of the points \(a, b\) . {[}To show that the condition is necessary, use Exercise 4 a) of Chapter I, § 11; to show that it is sufficient, let A (resp. B) be the set of all \(y \in X\) such that there is a connected subset of \(\mathbb{C}\{y\}\) which is closed in X and contains \(a\) and \(x\) (resp. \(b\) and \(x\) ); show that A and B are non-empty and open and disjoint from each other, and that \(\mathbb{C}\{x\} = A \cup B\) .

\begin{enumerate}
\def\labelenumi{\alph{enumi})}
\setcounter{enumi}{1}
\tightlist
\item
  Let X be a compact connected space which is irreducible between two of its points \(a, b\) (Chapter II, § 4, Exercise 19). Show that if for each pair of distinct points \(x, y\) of X there is a compact connected set in X which contains \(y\) and one of \(a, b\) but does not contain \(x\) , then there is a linear ordering on X such that \(\mathfrak{C}_{0}(\mathbf{X})\) is the given topology on X (Exercises 15 and 16).
\end{enumerate}

\begin{enumerate}
\def\labelenumi{\arabic{enumi})}
\setcounter{enumi}{19}
\item
  In the construction of Exercise 23 of Chapter I, § 9, take \(\mathbf{X}_0\) to be the interval {[}0, 1{]} of R. Show that the homeomorphisms of X onto itself are the same as the homeomorphisms of \(\mathbf{X}_0\) onto itself, although the topology of X is strictly finer than the topology of \(\mathbf{X}_0\) {[}use Exercise 20 b) of Chapter I, § 8{]}.
\item
  Show that if \(r > 2\) there is no \(r\) -fold transitive group of homeomorphisms of \(\mathbf{R}\) onto itself (consider the subgroup leaving two points of \(\mathbf{R}\) fixed).
\end{enumerate}

\BourbakiExerciseGroup

\begin{enumerate}
\def\labelenumi{\arabic{enumi})}
\tightlist
\item
  Let \(x\) be a real number \(\geqslant 0\) and let \(p\) and \(q\) be two integers \(> 0\) . Prove the relations
\end{enumerate}

\[
(x ^ {1 / p}) ^ {1 / q} = x ^ {1 / p q}, \quad (x ^ {p}) ^ {1 / q} = (x ^ {1 / q}) ^ {p}, \quad x ^ {1 / p} x ^ {1 / q} = (x ^ {p + q}) ^ {1 / p q}.
\]

\begin{enumerate}
\def\labelenumi{\arabic{enumi})}
\setcounter{enumi}{1}
\tightlist
\item
  \begin{enumerate}
  \def\labelenumii{\alph{enumii})}
  \tightlist
  \item
    On the polynomial ring \(\mathbf{A} = \mathbb{R}[\mathbf{X}]\) , consider the linearly ordered ring structure for which the elements \(\geqslant 0\) are \(0\) and all polynomials \(\neq 0\) whose leading coefficient is \(>0\) . Show that the topology \(\mathcal{G}_{0}(\mathbf{A})\) is not compatible with the ring structure of \(\mathbf{A}\) .
  \end{enumerate}
\end{enumerate}

\begin{enumerate}
\def\labelenumi{\alph{enumi})}
\setcounter{enumi}{1}
\item
  If K is an ordered field the topology \(\mathcal{G}_{0}(\mathbf{K})\) is compatible with the ring structure of K, and the bounded sets for this ring topology (Chapter III, § 6, Exercise 12) are the bounded sets for the order structure. The topology \(\mathcal{G}_{0}(\mathbf{K})\) is locally retrobounded (Chapter III, § 6, Exercise 22) and in particular is compatible with the field structure of K; moreover the completion \(\hat{\mathbf{K}}\) of K with respect to this topology, which is a field (loc. cit.) is canonically endowed with the structure of an ordered field {[}§ 1, Exercise 3 b){]}.
\item
  On the field \(\mathbf{K} = \mathbf{Q}(\sqrt{2})\) , considered as an ordered subfield of R, consider the topology C which is such that \((x, y) \to x + y\sqrt{2}\) is a homeomorphism of \(Q^{2}\) onto K, \(Q^{2}\) carrying the product topology. This topology C is finer than \(\mathcal{G}_{0}(\mathbf{K})\) and is compatible with the field structure of K; but the completion of this topological field is not a field.
\end{enumerate}

\(d\) ) The field \(\mathbf{K} = \mathbf{Q}(\sqrt[3]{2})\) , considered as a vector space over \(\mathbf{Q}\) , is the direct sum of the vector subspace \(\mathbf{G}'\) generated by \(\mathbf{I}\) and \(\sqrt[3]{2}\) , and the vector subspace \(\mathbf{G}''\) generated by \(\sqrt[3]{4}\) . Consider the topologies on \(\mathbf{G}'\) and \(\mathbf{G}''\) induced by the topology of \(\mathbf{R}\) ; let \(\mathfrak{C}\) be the topology on \(\mathbf{K}\) which is the product of the topologies of \(\mathbf{G}'\) and \(\mathbf{G}''\) . Show that \(\mathfrak{C}\) is compatible with the additive group structure of \(\mathbf{K}\) and is finer than \(\mathfrak{C}_0(\mathbf{K})\) ( \(\mathbf{K}\) being considered as an ordered subfield of \(\mathbf{R}\) ) but is not compatible with the ring structure of \(\mathbf{K}\) .

¶ 3) a) If f is an isomorphism of the field R onto a subfield of R show that f must be the identity mapping of R onto itself. {[}Remark that all rational numbers are invariant under f and that

\[
f (x ^ {2}) = (f (x)) ^ {2} \geqslant 0,
\]

so that \(f\) is increasing{]}.

\begin{enumerate}
\def\labelenumi{\alph{enumi})}
\setcounter{enumi}{1}
\tightlist
\item
  Let \(K_{0}\) be the field \(\mathbf{R}(\mathbf{X})\) of rational fractions in one indeterminate over R, ordered by taking the polynomials \textgreater0 to be those whose leading coefficient is \textgreater0. Let K be the maximal ordered algebraic extension of \(K_{0}\) . Show that there is an order-preserving automorphism f of K such that \(f(\xi)=\xi\) for all \(\xi\in\mathbb{R}\) and \(f(\mathbf{X})=\mathbf{X}^{2}(*)\) .
\end{enumerate}

\BourbakiExerciseGroup

\begin{enumerate}
\def\labelenumi{\arabic{enumi})}
\item
  Show that every bounded half-open interval of \(\mathbf{R}\) is homeomorphic to every unbounded closed interval of \(\mathbf{R}\) . If \(a < b\) , then no two of the three intervals \(]a, b[, [a, b[, [a, b]]\) are homeomorphic.
\item
  Show that the mapping \(x \to x / (1 + |x|)\) is a homeomorphism of \(\mathbf{R}\) onto the open interval \(] - 1, 1[\) .
\item
  Let I denote the interval \([0, 1]\) of \(\mathbf{R}\) , and let \(f\) be a homeomorphism of I onto itself such that \(f(0) = 0\) and \(f(1) = 1\) . Show that there is a continuous mapping \(g\) of \(\mathbf{I} \times \mathbf{I}\) into I such that: (i) for each \(x \in \mathbf{I}\) , \(g(x, 0) = x\) and \(g(x, 1) = f(x)\) ; (ii) for each \(y \in \mathbf{I}\) the partial mapping \(x \to g(x, y)\) is a homeomorphism of I onto itself such that \(g(0, y) = 0\) and \(g(1, y) = 1\) .
\item
  Show that the uniformity induced on \(\mathbf{R}\) by the unique uniformity of \(\overline{\mathbf{R}}\) is strictly coarser than the additive uniformity of \(\mathbf{R}\) .
\item
  As the point \((x, y)\) tends to \((+\infty, -\infty)\) while remaining in \(\mathbf{R} \times \mathbf{R}\) , shows that the set of cluster points of the function \(x + y\) is \(\overline{\mathbf{R}}\) .
\end{enumerate}

Likewise, as \((x,y)\) tends to \((0,+\infty)\) while remaining in \(R\times R\) , the set of cluster points of xy is \(\overline{R}\) .

\begin{enumerate}
\def\labelenumi{\arabic{enumi})}
\setcounter{enumi}{5}
\item
  Show that every rational function \(\mathbf{P}(x)/\mathbf{Q}(x)\) with coefficients in R, defined at all points x of R such that \(\mathbf{Q}(x)\neq0\) , can be extended by continuity to the points \(+\infty\) and \(-\infty\) (taking its values in \(\overline{R}\) ).
\item
  Let X be a linearly ordered set.
\end{enumerate}

\begin{enumerate}
\def\labelenumi{\alph{enumi})}
\item
  Let \(X_{1}\) be the completion of X (S et Theory, Chapter III, § 1, Exercise 15). \(X_{1}\) is a linearly ordered set whose elements are the subsets A of X such that; (i) if \(x \in A\) and \(y \leqslant x\) , then \(y \in A\) ; (ii) if A has a least upper bound in X, this bound belongs to A {[}which is equivalent to saying that A is closed in \(\mathcal{G}_{0}(X)\) {]}; the elements of \(X_{1}\) are ordered by inclusion. \(X_{1}\) is compact with respect to the topology \(\mathcal{G}_{0}(X_{1})\) {[}§ 2, Exercise 6 a){]}. The mapping \(x \to ]\leftarrow,x]\) of X into \(X_{1}\) is strictly order-preserving, and if we identify X with its image under this mapping, then X is dense in \(X_{1}\) and \(\mathcal{G}_{0}(X)\) is induced by \(\mathcal{G}_{0}(X_{1})\) . \(X_{1}\) is connected if and only if X is without gaps (§ 2, Exercise 7).
\item
  For each cardinal c, give an example of a linearly ordered set X which is connected in the topology \(\mathfrak{G}_{0}(\mathbf{X})\) and is such that every fundamental system of neighbourhoods of any point of X has cardinal \(\geqslant c\) {[}use a) and the method of §1, Exercise 4 b){]}.
\item
  With the notation of \(a\) , let \(\mathbf{X}_{2}\) be the subset of the lexicographic product \(\mathbf{X}_{1} \times \left\{-\mathbf{i}, 0, \mathbf{i}\right\}\) which is the complement of the set consisting of: (i) the points \((x, 0)\) where \(x \notin \mathbf{X}\) ; (ii) the points \((x, 1)\) where \(x \in \mathbf{X}\) and the set of all \(y > x\) in \(\mathbf{X}\) has a smallest element; (iii) the points \((x, -1)\) where \(x \in \mathbf{X}\) and the set of all \(y < x\) in \(\mathbf{X}\) has a greatest element. Show that \(\mathbf{X}_{2}\) , endowed with the topology \(\mathfrak{C}_{0}(\mathbf{X}_{2})\) , is compact and totally disconnected, that \(\mathbf{X}\) can be identified, by means of the strictly order-preserving mapping \(x \to (x, 0)\) , with a dense subset of \(\mathbf{X}_{2}\) , and that the topology induced by \(\mathfrak{C}_{0}(\mathbf{X}_{2})\) on \(\mathbf{X}\) is discrete.
\item
  Let \(X'\) be a linearly ordered set containing X which induces the given order structure on X, is compact in the topology \(\mathcal{C}_{0}(\mathbf{X}')\) and is such that X is dense in \(X'\) with respect to this topology. Show that there is a continuous order-preserving surjective mapping \(f: X' \to X_{1}\) and a continuous order-preserving surjective mapping \(g: X_{2} \to X'\) , both of which induce the identity mapping on X.
\end{enumerate}

\BourbakiExerciseGroup

\begin{enumerate}
\def\labelenumi{\arabic{enumi})}
\item
  Let \(X_{1}, X_{2}\) be two directed sets and let f be a real-valued function defined on \(X_{1} \times X_{2}\) such that, for each \(x_{1} \in X_{1}\) , the mapping \(x_{2} \to f(x_{1}, x_{2})\) is increasing on \(X_{2}\) , and for each \(x_{2} \in X_{2}\) , the mapping \(x_{1} \rightarrow f(x_{1}, x_{2})\) is increasing on \(X_{1}\) . Show that f has a limit in \(\overline{R}\) with respect to the product of the section filters of \(X_{1}\) and \(X_{2}\) , and that this limit is the least upper bound of f.
\item
  \begin{enumerate}
  \def\labelenumii{\alph{enumii})}
  \tightlist
  \item
    Let \(f\) be a real-valued function defined on a set \(\mathbf{X}\) , let \(\mathbf{A}\) be a non-empty subset of \(\mathbf{X}\) , and let \(\varphi\) be an increasing real-valued function defined on \(f(\mathbf{A})\) . If \(\varphi\) is continuous at the point \(a = \sup_{x \in \mathbf{A}} f(x)\) , then \(\sup_{x \in \mathbf{A}} \varphi(f(x)) = \varphi(\sup_{x \in \mathbf{A}} f(x))\) .
  \end{enumerate}
\end{enumerate}

\begin{enumerate}
\def\labelenumi{\alph{enumi})}
\setcounter{enumi}{1}
\tightlist
\item
  Let f be a real-valued function defined on a set X filtered by a filter F, and let \(\varphi\) be a real-valued function defined on an open neighbourhood V of the set of cluster points of f with respect to F. Show that if \(\varphi\) is increasing and continuous on V, then
\end{enumerate}

\[
\lim \sup _ {\mathfrak {F}} (\varphi \circ f) = \varphi (\lim \sup _ {\mathfrak {F}} f).
\]

\begin{enumerate}
\def\labelenumi{\arabic{enumi})}
\setcounter{enumi}{2}
\item
  Let \(f\) be a real-valued function defined on an infinite set \(\mathbf{X}\) and let \(\mathfrak{G}\) be the filter of complements of finite subsets of \(\mathbf{X}\) . Show that \(\limsup_{\mathfrak{G}} f\) is the least upper bound of the set of real numbers \(x\) such that the set \(f([x, +\infty])\) is infinite.
\item
  Let \(f, g\) be two real-valued functions defined on a set \(X\) filtered by a filter \(\mathfrak{F}\) . Show that
\end{enumerate}

\[
\lim \sup _ {\mathfrak {F}} (\sup (f, g)) = \sup (\lim \sup _ {\mathfrak {F}} f, \lim \sup _ {\mathfrak {F}} g).
\]

Give an example of a set X filtered by a filter F, and an infinite family \((f_{i})\) of real-valued functions defined on X, such that

\[
\sup _ {t} (\lim \sup _ {\mathfrak {F}} f _ {t}) <   \lim \sup _ {\mathfrak {F}} (\sup _ {t} f _ {t}).
\]

\begin{enumerate}
\def\labelenumi{\arabic{enumi})}
\setcounter{enumi}{4}
\tightlist
\item
  Let \(f, g\) be two real-valued functions defined on a filtered set. Show that if \(\lim g\) exists and is \(\geqslant 0\) , then
\end{enumerate}

\[
\lim \sup f g = (\lim \sup f) (\lim g)
\]

whenever both sides are defined.

\begin{enumerate}
\def\labelenumi{\arabic{enumi})}
\setcounter{enumi}{5}
\item
  Let \(X_{1}\) (resp. \(X_{2}\) ) be a set filtered by a filter \(F_{1}\) (resp. \(F_{2}\) ) and let f be a real-valued function defined on \(X_{1} \times X_{2}\) . Show by an example that the three numbers \(\limsup_{\mathfrak{F}_{i} \times \mathfrak{F}_{2}} f(x_{1}, x_{2})\) , \(\limsup_{\mathfrak{F}_{i}} (\limsup_{\mathfrak{F}_{2}} f(x_{1}, x_{2}))\) , \(\limsup_{\mathfrak{F}_{i}} (\limsup_{\mathfrak{F}_{i}} f(x_{1}, x_{2}))\) are in general distinct.
\item
  Let \(f\) be a real-valued function defined on a closed subset \(A\) of \(\overline{\mathbf{R}}\) . Let \(\limsup_{x \to a, x \geqslant a} f(x)\) {[}resp. \(\limsup_{x \to a, x \leqslant a} f(x)\) {]} denote the upper limit of \(f(x)\) as \(x\) tends to a point \(a \in A\) while remaining in \(A \cap [a, +\infty]\) .
\end{enumerate}


(resp. A ∩ {[}---∞, a{]}). Show that the set of points \(a \in A\) such that \(\limsup_{x \to a, x \geqslant a} f(x) \neq \limsup_{x \to a, x \leqslant a} f(x)\) is countable. {[}Prove that, for each pair of rational numbers p, q such that p \textless{} q, the set \(G_{p,q}\) of points \(a \in A\) such that

\[
\lim _ {x \rightarrow a, x \geqslant a} \sup f (x) \leqslant p <   q \leqslant \lim _ {x \rightarrow a, x \leqslant a} \sup f (x)
\]

is countable, by using Exercise 1 of § 2{]}.

\begin{enumerate}
\def\labelenumi{\arabic{enumi})}
\setcounter{enumi}{7}
\item
  Deduce from Exercise 7 that a real-valued function \(f\) defined on a closed subset \(\mathbf{A}\) of \(\overline{\mathbf{R}}\) , and monotone on \(\mathbf{A}\) , is continuous on \(\mathbf{A}\) except at the points of a countable subset of \(\mathbf{A}\) .
\item
  Let X be a topological space with a countable base and let f be a real-valued function defined on X. Show that the set of points \(x \in X\) , such that \(\lim_{x \to a, x \neq a} f(x)\) exists and differs from \(f(a)\) , is countable (same method as in Exercise 7).
\item
  Let X be a topological space which has a countable base \((\mathbf{U}_{n})\) and let f be a real-valued function defined on X; f is said to attain a strict relative maximum at a point \(a \in X\) if there is a neighbourhood V of a such that \(f(x) < f(a)\) for each \(x \in V\) other than x = a. Show that the set M of points of X, at which f attains a strict relative maximum, is countable. (Consider the set of \(U_{n}\) such that f attains, at some point of \(U_{n}\) , a strict relative maximum equal to its least upper bound in \(U_{n}\) , and show that there is a mapping of this set onto M).
\item
  Let \((u_n)\) be a sequence of real numbers \(>0\) such that \(\lim_{n\to \infty}u_n = 0\) . Show that there is an infinite number of indices \(n\) such that \(u_{n}\geqslant u_{m}\) for all \(m\geqslant n\) .
\item
  Let \((u_n)\) be a sequence of real numbers \(>0\) such that
\end{enumerate}

\[
\lim _ {n \to \infty} \inf u _ {n} = 0.
\]

Show that there is an infinite number of indices \(n\) such that \(u_{n} \leqslant u_{m}\) for all \(m \leqslant n\) .

\begin{enumerate}
\def\labelenumi{\arabic{enumi})}
\setcounter{enumi}{12}
\item
  Let \((u_{n})\) be a sequence of finite real numbers and let \((\varepsilon_{n})\) be a sequence of numbers \(\geqslant 0\) , such that \(\lim_{n\to \infty}\varepsilon_n = 0\) and \(u_{n + 1}\geqslant u_n - \varepsilon_n\) for all integers \(n\geqslant 0\) . Let \(a = \liminf_{n\to \infty}u_n\) and let \(b = \limsup_{n\to \infty}u_n\) ; show that the set of cluster points of the sequence \((u_{n})\) is the interval \([a,b]\) .
\item
  Let \((r_n)\) be an increasing sequence of finite real numbers \(>0\) , such that \(\lim r_n = +\infty\) . For each finite real number \(r > 0\) , let \(\mathbf{N}(r)\) denote the largest index \(n\) such that \(r_n \leqslant r\) . Show that
\end{enumerate}

\[
\limsup _ {r \to \infty} \frac {\mathrm{N} (r)}{r} = \limsup _ {n \to \infty} \frac {n}{r _ {n}}, \quad \liminf _ {r \to \infty} \frac {\mathrm{N} (r)}{r} = \liminf _ {n \to \infty} \frac {n}{r _ {n}}.
\]

¶ 15) Let \((x_n)\) be a sequence of finite real numbers and let \((p_n)\) be a sequence of finite real numbers \(\geqslant 0\) such that \(\lim_{n\to \infty}\left(\sum_{i=0}^{n}p_i\right) = +\infty\) . Let \(y_n = \left(\sum_{i=0}^{n}p_ix_i\right) / \left(\sum_{i=0}^{n}p_i\right)\) for those values of \(n\) such that \(\sum_{i=1}^{\infty}p_i \neq 0\) . Show that \(\liminf_{n\to \infty}x_n \leqslant \liminf_{n\to \infty}y_n \leqslant \limsup_{n\to \infty}y_n \leqslant \limsup_{n\to \infty}x_n\) .

Let H be any non-empty subset of the set of cluster points of the sequence \((x_{n})\) in \(\overline{R}\) . Show that the sequence \((p_{n})\) of finite real numbers \(\geqslant0\) can be determined so that \(\lim_{n\to\infty}\left(\sum_{i=0}^{n}p_{i}\right)=+\infty\) and so that the set of cluster points of the corresponding sequence \((y_{n})\) contains H. {[}Reduce to the case where H is countable, then define \((p_{n})\) by induction, taking its terms to be o or 1{]}.

Deduce that the sequence \((x_{n})\) converges in \(\overline{\mathbf{R}}\) if and only if the sequence \((y_{n})\) converges in \(\overline{\mathbf{R}}\) for every sequence \((p_n)\) of numbers \(\geqslant 0\) such that \(\lim_{n\to \infty}\left(\sum_{i = 0}^{n}p_i\right) = +\infty .\)

\begin{enumerate}
\def\labelenumi{\arabic{enumi})}
\setcounter{enumi}{15}
\tightlist
\item
  Let \(x_0, y_0\) be two real numbers such that \(0 < y_0 \leqslant x_0\) . For \(n > 0\) we define two sequences \((x_n), (y_n)\) inductively by the relations
\end{enumerate}

\[
x _ {n + 1} = \left(x _ {n} + y _ {n}\right) / 2, \quad y _ {n + 1} = \sqrt {x _ {n} y _ {n}}.
\]

Show that the two sequences \((x_{n})\) , \((y_{n})\) tend to the same limit \(\alpha\) (the ``arithmetic-geometric mean'' of \(x_{0}, y_{0}\) ); also there exist numbers \(a > 0\) and \(\gamma\) such that \(0 < \gamma < 1\) and \(x_{n} - y_{n} \leqslant a\gamma^{2n}\) for each \(n\) {[}observe that \(x_{n+1} - y_{n+1} = (x_{n} - y_{n})^{2}/4(x_{n+1} + y_{n+1})\) {]}.

\begin{enumerate}
\def\labelenumi{\arabic{enumi})}
\setcounter{enumi}{16}
\tightlist
\item
  Let \(g\) be a mapping of \(]0, i]\) into \([-1, 1]\) such that
\end{enumerate}

\[
\lim _ {x \rightarrow 0, x > 0} g (x) = 0.
\]

Show that there is a continuous increasing mapping \(g_{2}\) and a continuous decreasing mapping \(g_{1}\) of {[}0, 1{]} into {[}-1, 1{]} such that \(g_{1}(0) = g_{2}(0) = 0\) .

and \(g_{1}(x) \leqslant g(x) \leqslant g_{2}(x)\) for \(0 < x \leqslant 1\) . {[}For each integer \(n > 0\) consider the greatest lower bound \(x_{n}\) of the numbers \(x\) such that \(g(x) \geqslant 1 / n\) {]}.

\begin{enumerate}
\def\labelenumi{\arabic{enumi})}
\setcounter{enumi}{17}
\tightlist
\item
  Extend the definitions and results of no. 1 to no. 6 to functions which take their values in a linearly ordered set X which is compact in the topology \(\mathfrak{G}_{0}(\mathbf{X})\) (cf.~§ 2, Exercise 6).
\end{enumerate}

\BourbakiExerciseGroup

\begin{enumerate}
\def\labelenumi{\arabic{enumi})}
\item
  Let \(f\) be a continuous mapping of an open interval \(\mathbf{I} \subset \mathbb{R}\) into \(\mathbb{R}\) . Show that if \(f(\mathbf{I})\) is open and if, for each \(y \in \mathbb{R}\) , the set \(\overline{f}^1(y)\) has at most two distinct points, then \(f\) is monotone.
\item
  Let B be a base of R considered as a vector space over the field Q (``Hamel base''). B is uncountable (§ 2, Exercise 2). Let φ be a bijection of a subset C ≠ B of B onto the set B. Define a mapping f of R into R as follows: if \(x = \sum_{\xi \in B} \lambda(\xi) \xi\) , where \(\lambda(\xi) \in \mathbf{Q}\) and \(\lambda(\xi) = 0\) for all but a finite number of elements \(\xi \in B\) , then \(f(x) = \sum_{\xi \in C} \lambda(\xi) \varphi(\xi)\) . Show that \(f(x + y) = f(x) + f(y)\) , but that, for each \(z \in R\) , \(f(z)\) is a dense subset of R, which implies that f is not bounded above nor below in any interval of R.
\item
  \begin{enumerate}
  \def\labelenumii{\alph{enumii})}
  \tightlist
  \item
    For each interval \(\mathbf{I} \subset \mathbb{R}\) , let \(\mathbf{G}(\mathbf{I})\) denote the group of all homeomorphisms of \(\mathbf{I}\) onto itself. Show that if \(\mathbf{I}\) and \(\mathbf{J}\) are any two intervals of \(\mathbb{R}\) , each containing more than one point, then \(\mathbf{G}(\mathbf{I})\) and \(\mathbf{G}(\mathbf{J})\) are isomorphic. In \(\mathbf{G}(\mathbf{I})\) , the set \(\mathbf{F}(\mathbf{I})\) of increasing homeomorphisms is a normal subgroup of index 2.
  \end{enumerate}
\end{enumerate}

\begin{enumerate}
\def\labelenumi{\alph{enumi})}
\setcounter{enumi}{1}
\item
  Let \(\mathbf{G} = \mathbf{G}(\mathbf{R})\) and \(\mathbf{F} = \mathbf{F}(\mathbf{R})\) . For each \(f \in \mathbf{G}\) and each \(x \in \mathbf{R}\) , let \(\sigma(x; f)\) denote \(\operatorname{sgn}(f(x) - x)\) . Let \(\mathbf{T}^{+}\) (resp. \(\mathbf{T}^{-}\) ) be the set of all \(f \in \mathbf{F}\) such that \(\sigma(x; f)\) is constant and equal to \(\mathfrak{I}\) (resp. \(-\mathfrak{I}\) ) on \(\mathbf{R}\) . Show that every element of \(\mathbf{T}^{+}\) (resp. \(\mathbf{T}^{-}\) ) is conjugate in \(\mathbf{F}\) to the translation \(x \to x + \mathfrak{I}\) (resp. \(x \to x - \mathfrak{I}\) ). {[}If \(f \in \mathbf{T}^{+}\) , consider the sequence of points \(f^{n}(0), n \in \mathbf{Z}\) {]}.
\item
  Let \(T = T^{+} \cup T^{-}\) . Show that every \(f \in F\) which does not belong to T is the product in F of two elements of \(T_{*}\) {[}Write \(f = f_{1}f_{2}\) in F, where \(f_{1}(x) = \sup (x, f(x))\) and \(f_{2}(x) = \inf (x, f(x))\) ; if g is the translation \(x \to x + 1\) , then \(f_{1}g\) and \(g^{-1}f_{2}\) belong to T{]}.
\item
  Two elements \(f, g\) are conjugate in \(F\) if and only if there exists \(s \in F\) such that \(\sigma(x; f) = \sigma(s(x); g)\) for all \(x \in \mathbb{R}\) . {[}Consider the components \(I_n\) of the open set of all \(x\) such that \(\sigma(x; f) \neq 0\) ; observe that \(F(I_n)\) and \(F\) are isomorphic (by \(a\) ) and use \(b\) {]}. In particular, if \([a, b]\) is an interval of \(R\) such that \(\sigma(a; f) = \sigma(b; f) = 0\) for some \(f \in F\) , then the element \(f'\) of \(F\) which coincides with \(f\) for \(x \leqslant a\) or \(x \geqslant b\) and is such that \(f'(x) = x + \varepsilon(f(x) - x)\) for \(a < x < b\) , where \(0 < \varepsilon < 1\) , is conjugate to \(f\) in \(F\) .
\item
  Let \(f \in F\) and let \(a, b, c, d\) be four real numbers such that \(a < c < b < d\) and \(f(x) - x = 0\) for \(x = a, x = b, x = c, x = d, \sigma(x; f) = 1\) for \(a < x < c\) and \(b < x < d\) . Let \(a', b'\) be two numbers such that \(a \leqslant a' < c\) and \(b < b' \leqslant d\) ; let \(J = [a, b], J' = [a', b']\) , and let \(f_1\) denote the restriction of \(f\) to \(J\) ; thus \(f_1 \in F(J)\) . Show that there is an increasing homeomorphism \(s\) of \(J\) onto \(J'\) such that the element \(f_2 = sf_1s^{-1}\) of \(F(J')\) is such that \(f_2(x) > f^{-1}(x)\) whenever \(a' < x < b'\) {[}use \(d\) {]}.
\end{enumerate}

\(f)\) Let \(f \in \mathbf{F}\) and let \([a, b]\) be an interval of \(\mathbf{R}\) such that \(f(x) - x = 0\) for \(x = a\) and \(x = b\) , \(\sigma(x; f) = -1\) for \(x < a\) and \(\sigma(x; f) = +1\) for \(x > b\) . Show that there exists \(g \in \mathbf{F}\) , conjugate to \(f\) in \(\mathbf{F}\) , such that \(g(x) > f(x)\) for all \(x \in \mathbf{R}\) (same method).

\(g\) ) Let \(\mathbf{H}^{+}\) (resp. \(\mathbf{H}^{-}\) ) denote the normal subgroup of \(\mathbf{F}\) consisting of all \(f\in \mathbf{F}\) such that \(f(x) = x\) in a neighbourhood of \(+\infty\) (resp. \(-\infty\) ). Show that \(\mathbf{H}^{+},\mathbf{H}^{-}\) and \(\mathbf{H} = \mathbf{H}^{+}\cap \mathbf{H}^{-}\) are the only normal subgroups of \(\mathbf{F}\) other than \(\mathbf{F}\) and \(\{e\}\) . {[}If \(\mathbf{N}\) is a normal subgroup of \(\mathbf{F}\) , other than \(\mathbf{F}\) , and if \(f\in \mathbf{N}\) does not belong to \(\mathbf{H}^{+}\) , show that there is an element \(g\in \mathbf{N}\) such that the set of all \(x\in \mathbf{R}\) for which \(\sigma (x;g) = +1\) is an interval \([a, + \infty ]\) , and that we have \(g(x) = x\) for \(x\leqslant a\) . To do this use the constructions of \(e)\) and \(f)\) , as well as \(b)\) and \(c)\) . To show that \(\mathbf{N}\) contains no subgroup, other than \(\mathbf{H}\) and \(\{e\}\) , which is normal in \(\mathbf{F}\) , consider an element \(f\in \mathbf{H}\) distinct from \(e\) ; let \(a\) and \(b\) be the greatest lower bound and least upper bound (both finite by hypothesis) of the set of \(x\in \mathbf{R}\) such that \(\sigma (x;f)\neq 0\) . Observe then that \(\mathrm{F}([a,b])\) is isomorphic to \(\mathbf{F}\) and use the preceding result{]}.

\begin{enumerate}
\def\labelenumi{\alph{enumi})}
\setcounter{enumi}{7}
\tightlist
\item
  Show that the group H is simple. (Same method as for proving that H contains no non-trivial subgroups normal in F).
\end{enumerate}

\begin{enumerate}
\def\labelenumi{\arabic{enumi})}
\setcounter{enumi}{3}
\item
  Let \(f\) be a lower semi-continuous function on a topological space \(\mathbf{X}\) , and let \(\varphi\) be an increasing lower semi-continuous function on \(f(\mathbf{E}) \subset \overline{\mathbf{R}}\) . Show that \(\varphi \circ f\) is lower semi-continuous on \(\mathbf{X}\) .
\item
  Let \(f\) be a lower semi-continuous function on a topological space \(\mathbf{X}\) . Show that if \(\mathbf{A}\) is any non-empty subset of \(\mathbf{X}\) , then \(\sup f(\overline{\mathbf{A}}) = \sup f(\mathbf{A})\) .
\item
  Let \(f\) be a real-valued function defined on a topological space \(X\) . The number (finite or infinite)
\end{enumerate}

\[
\omega (a; f) = \lim _ {x \to a} \sup f (x) - \lim _ {x \to a} \inf f (x)
\]

is called the oscillation of \(f\) at \(a \in \mathbf{X}\) , whenever the right-hand side is defined.

\begin{enumerate}
\def\labelenumi{\alph{enumi})}
\item
  Show that \(x \to \omega(x; f)\) is upper semi-continuous on the subspace \(A\) of \(X\) where this mapping is defined.
\item
  If \(f\) is finite on \(\mathbf{X}\) , then \(\mathbf{A} = \mathbf{X}\) and for each \(a \in \mathbf{X}\) we have
\end{enumerate}

\[
\omega (a; f) = \lim _ {(x, y) \rightarrow (a, a)} \sup (f (x) - f (y)).
\]

\begin{enumerate}
\def\labelenumi{\alph{enumi})}
\setcounter{enumi}{2}
\tightlist
\item
  Let \(f\) be a lower semi-continuous finite real-valued function on \(\mathbf{X}\) . Show that if, at a point \(a \in \mathbf{X}, \omega(a; f)\) has a finite value, then
\end{enumerate}

\[
\liminf _ {x \to a} \omega (x; f) = 0.
\]

{[}Assuming the result false, show that there would exist points x arbitrarily near to a such that \(f(x)\) is as large as we please{]}.

\begin{enumerate}
\def\labelenumi{\alph{enumi})}
\setcounter{enumi}{3}
\tightlist
\item
  For each rational number \(r = p / q\) in irreducible form (with \(q > 0\) ), put \(f(r) = q\) . Show that \(f\) is lower semi-continuous on \(\mathbf{Q}\) , but that at each point \(r \in \mathbf{Q}\) we have \(\omega(r; f) = +\infty\) .
\end{enumerate}

¶ 7) Let X be a locally compact space and f a lower semi-continuous function on X.

\begin{enumerate}
\def\labelenumi{\alph{enumi})}
\tightlist
\item
  If \(\limsup_{x\to a}f(x)=+\infty\) for every \(a\in\mathbf{X}\) , show that the set \(\overline{f}(+\infty)\) is dense in \(\mathbf{X}\) {[}for each \(a\in\mathbf{X}\) and each neighbourhood \(\mathbf{V}\) of \(a\) , define a sequence \((\mathbf{U}_{n})\) of relatively compact open sets such that \(\overline{\mathbf{U}}_{n+1}\subset\mathbf{U}_{n}\) and \(f(x)>n\) for all \(x\in\mathbf{U}_{n}\) {]}.
\end{enumerate}

* b) Let \(n \to r_n\) be a bijection of \(\mathbf{N}\) onto the set of rational numbers belonging to {[}0, 1{]}, and let \(\varphi(x) = \mathrm{i}/\sqrt{|x|}\) {[}with \(\varphi(0) = +\infty\) {]}; then the function \(f(x) = \sum_{n=0}^{\infty} 2^{-n} \varphi(x - r_n)\) is lower semi-continuous on {[}0, 1{]}, \(\overline{f}(+\infty)\) is dense in this interval, and so is its complement {[}to prove the last point, observe that the series is convergent whose general term is

\[
2 ^ {- n} \int_ {0} ^ {1} \varphi (x - r _ {n}) d x. *
\]

(Integration, Chapter IV, § 3, no. 6, Theorem 5).{]} *

\begin{enumerate}
\def\labelenumi{\alph{enumi})}
\setcounter{enumi}{2}
\item
  If \(f\) is finite, show that the set of points \(x\) such that \(\omega(x; f)\) is finite is dense {[}use a){]}.
\item
  Show that the set of points of X at which f is continuous (f being finite or not) is dense. {[}Reduce to the case in which f is bounded, by replacing f by ( f/(1 + \textbar f\textbar) ) (Exercise 4); observe that \(1/\omega(x; f)\) is lower semi-continuous and use a) and Exercise 6c){]}.
\end{enumerate}

\begin{enumerate}
\def\labelenumi{\arabic{enumi})}
\setcounter{enumi}{7}
\item
  Let X be a topological space, and let A be a closed subset of \(X \times \overline{R}\) . Show that the mapping \(x \to \inf(A(x))\) of \(pr_{1}A\) into \(\overline{R}\) is lower semi-continuous. Conversely, if \(f: X \to \overline{R}\) is a lower semi-continuous function, then the subset B of \(X \times \overline{R}\) consisting of pairs \((x, y)\) such that \(y \geqslant f(x)\) is closed in \(X \times \overline{R}\) .
\item
  \begin{enumerate}
  \def\labelenumii{\alph{enumii})}
  \tightlist
  \item
    Let X and Y be two Hausdorff spaces and let \(\pi: X \to Y\) be a proper mapping (Chapter I, § 10, no. 1). Let \(g\) be a lower semi-continuous real-valued function on X. For each \(y \in Y\) , let \(f(y)\) be the greatest lower bound of \(g\) in the set \(\overline{\pi}^{1}(y)\) {[}thus \(f(y) = +\infty\) if \(\overline{\pi}^{1}(y) = \emptyset\) {]}. Show that \(f\) is lower semi-continuous on Y. {[}Use the fact that \(\pi(X)\) is closed, that \(\overline{\pi}^{1}(y)\) is compact for all \(y \in \pi(X)\) , and that every neighbourhood of \(\overline{\pi}^{1}(y)\) contains a neighbourhood which is saturated with respect to the equivalence relation \(\pi(x) = \pi(x')\) .{]}
  \end{enumerate}
\end{enumerate}

\begin{enumerate}
\def\labelenumi{\alph{enumi})}
\setcounter{enumi}{1}
\tightlist
\item
  Let \(g\) be the continuous function \(|x_1 x_2 - 1|\) defined on \(\mathbf{R} \times ]0, +\infty[\) , and for each \(x_1 \in \mathbb{R}\) let \(f(x_1) = \inf_{x_1 > 0} g(x_1, x_2)\) . Show that \(f\) is not lower semi-continuous.
\end{enumerate}

\begin{enumerate}
\def\labelenumi{\arabic{enumi})}
\setcounter{enumi}{9}
\tightlist
\item
  Extend Definition 1, Proposition 1 and Theorem 3 to functions defined on a topological space X which take their values in an arbitrary linearly ordered set Y. If f and g are two mappings of X into Y, both lower semi-continuous at a point a, then inf (f, g) and sup (f, g) are lower semi-continuous at a. So is \(f + g\) if Y is a linearly ordered commutative group.
\end{enumerate}

If \(\mathbf{Y}\) is compact in the topology \(\mathfrak{G}_{0}(\mathbf{Y})\) ( \(\S 2\) , Exercise 6), extend Theorem 4 and Propositions 3 and 4 to functions with values in \(\mathbf{Y}\) .

¶ 11) Let \(\overline{\mathbf{R}}\) carry the right topology (Chapter I, § 1, Exercise 2). A mapping \(f\) of a topological space \(\mathbf{X}\) into \(\overline{\mathbf{R}}\) is continuous in this topology if and only if \(f\) has a relative minimum at every point of \(\mathbf{X}\) . If \(\mathbf{X} = \mathbf{R}\) (with the usual topology), show that the set \(f(\mathbf{X})\) is then countable. {[}Suppose that the result is false, and that the interval \([\alpha, \beta] \cap f(\mathbf{X})\) is uncountable; for each \(y \in [\alpha, \beta]\) , let \(\mathbf{U}(y) = \overline{f}([y, +\infty])\) ; \(\mathbf{U}(y)\) is an open subset in \(\mathbf{R}\) . Let \((\mathbf{I}_n)\) be the sequence (finite or infinite) of components of \(\mathbf{U}(\beta)\) , and for each \(n\) let \(x_n\) be a point of \(\mathbf{I}_n\) ; for each \(y \in [\alpha, \beta]\) , let \(g_n(y)\) (resp. \(h_n(y) \) ) be the left-hand (resp. right-hand) end-point of the component of \(\mathbf{U}(y)\) which contains \(x_n\) . Show that for at least one value of \(n\) one of the functions \(g_n, h_n\) must take an uncountable infinity of distinct values in \([\alpha, \beta]\) ; now use Exercise 8 of § 5{]}.

\begin{enumerate}
\def\labelenumi{\arabic{enumi})}
\setcounter{enumi}{11}
\item
  Let \(f\) be a continuous non-constant finite real-valued function on a compact interval \([a, b]\) of \(\mathbf{R}\) , such that \(f(a) = f(b) = 0\) ; let \(Z\) be the closed subset of \([a, b]\) consisting of points \(x\) for which \(f(x) = 0\) , and let \(c\) be the largest of the lengths of the intervals contiguous to \(Z\) in \([a, b]\) . Show that, for each \(t\) such that \(0 < t < c\) , there is a point \(x \in [a, b]\) such that \(x + t \in [a, b]\) and \(f(x + t) = f(x)\) .
\item
  Let \(f\) be a continuous finite real-valued function on a compact interval \([a, b]\) of \(\mathbf{R}\) , and suppose that \(f\) is not strictly monotone. Show that there exists a point \(x_0 \in ]a, b[\) such that, for each \(\varepsilon > 0\) , there are two points \(y, z\) in \([a, b]\) such that \(x_0 - \varepsilon < y < x_0 < z < x_0 + \varepsilon\) and \(f(y) = f(z)\) .
\item
  Let \(f\) be a continuous mapping of \(\mathbf{R}\) into itself.
\end{enumerate}

\begin{enumerate}
\def\labelenumi{\alph{enumi})}
\tightlist
\item
  Show that if \(f\) is uniformly continuous on \(\mathbf{R}\) , then there exist two real numbers \(\alpha \geqslant 0\) and \(\beta \geqslant 0\) such that \(|f(x)| \leqslant \alpha |x| + \beta\) for all \(x \in \mathbf{R}\) . b) Show that if \(f\) is monotone and bounded in \(\mathbf{R}\) , then \(f\) is uniformly continuous on \(\mathbf{R}\) .
\end{enumerate}

¶ 15) Let X be the Alexandroff half-line {[}§ 2, Exercise 12 b){]}, and let X' be its compactification by adjunction of a point at infinity ω. Show that, for every continuous mapping f: X → X, either we have

\(\lim_{x \to \omega, x \in \mathbf{X}} f(x) = \omega,\) or else there exists \(z \in \mathbf{X}\) such that \(f\) is constant

in the interval \([z, \rightarrow[\) . {[}Show that if, as \(x\) tends to \(\omega, f\) has a cluster point \(c \in \mathbf{X}\) , then we must have \(\lim_{x \to \omega, x \in \mathbf{X}} = c\) ; prove that \(f\) can have no other cluster point in X', by using the fact that every increasing sequence converges in X. Then use Exercise 12 b) of § 2 and the fact that in X' every countable intersection of neighbourhoods of ω is a neighbourhood of ω{]}.

\BourbakiExerciseGroup

\begin{enumerate}
\def\labelenumi{\arabic{enumi})}
\tightlist
\item
  Let \((x_{t})\) be a family of real numbers, all of which belong to the interval \(A' = ] - \infty, +\infty]\) (resp. \(A'' = [-\infty, +\infty[\) ). The family \((x_{t})\) is summable in \(\overline{\mathbf{R}}\) if and only if one of the following conditions is satisfied:
\end{enumerate}

\begin{enumerate}
\def\labelenumi{\alph{enumi})}
\item
  At least one of the numbers \(x_{t}\) is equal to \(+\infty\) (resp. \(-\infty\) ).
\item
  At least one of the two sums \(\sum_{i} x_{i}^{+}, \sum_{i} x_{i}^{-}\) is finite.
\end{enumerate}

In the first case we have \(\sum_{i} x_{i} = +\infty\) (resp. \(-\infty\) ); in the second case, \(\sum_{i} x_{i} = \sum_{i} x_{i}^{+} - \sum_{i} x_{i}^{-}\) .

\begin{enumerate}
\def\labelenumi{\arabic{enumi})}
\setcounter{enumi}{1}
\item
  Let \((x_{i})\) be a family of real numbers all of which belong to the interval \(\mathbf{A}'\) (resp. \(\mathbf{A}''\) ) and which satisfies condition \(b\) ) of Exercise 1. Show that every subfamily of \((x_{i})\) is summable in \(\overline{\mathbf{R}}\) , and that the sum of the \(x'\) is associative (cf.~Chapter III, § 5, no. 3, Theorem 2).
\item
  Let \((x_{n})_{n\geqslant 0}\) be a decreasing sequence of finite real numbers \(\geqslant 0\) . This sequence is summable in \(\mathbf{R}\) if and only if the sequence \((2^{n}x_{2^{n}})_{n\geqslant 0}\) is summable in \(\mathbf{R}\) (``Cauchy's condensation test''). * Deduce that the series whose general term is \(1 / n^{\alpha}\) is convergent if \(\alpha >1\) and not convergent if \(0 < \alpha \leqslant 1\) . *
\item
  Let \((a_{n})\) be any sequence of finite real numbers \(\geqslant 0\) . Show that the sequence \(\left(\frac{a_n}{(1 + a_1)(1 + a_2)\ldots(1 + a_n)}\right)\) is summable in \(\mathbf{R}\) (express each term of this sequence as a difference).
\item
  Let \((d_n)\) be a sequence of finite real numbers \(\geqslant 0\) such that
\end{enumerate}

\[
\sum_ {n = 0} ^ {\infty} d _ {n} = + \infty .
\]

What can be said about the convergence of the series whose general terms are

\[
\frac {d _ {n}}{\mathrm{I} + d _ {n}}, \quad \frac {d _ {n}}{\mathrm{I} + n d _ {n}}, \quad \frac {d _ {n}}{\mathrm{I} + n ^ {2} d _ {n}}, \quad \frac {d _ {n}}{\mathrm{I} + d _ {n}}?
\]

\begin{enumerate}
\def\labelenumi{\arabic{enumi})}
\setcounter{enumi}{5}
\item
  Prove that \(s = \sum_{n=1}^{\infty} (1/n)\) is equal to \(+\infty\) by showing that \(s \geqslant s + \frac{1}{2}\) (find lower bounds, as functions of \(s\) , for the sum of the terms with even indices and the sum of the terms with odd indices).
\item
  Show that, for each integer \(p > \mathbf{r}\) ,
\end{enumerate}

\[
\sum_ {n = 1} ^ {\infty} \frac {(- \mathrm{I}) ^ {n - 1}}{n ^ {p}} = (\mathrm{I} - 2 ^ {1 - p}) \sum_ {n = 1} ^ {\infty} \frac {\mathrm{I}}{n ^ {p}}.
\]

\begin{enumerate}
\def\labelenumi{\arabic{enumi})}
\setcounter{enumi}{7}
\item
  Show that, if \(m\) is an integer \(> 0\) , then \(\sum_{n \geqslant 1, n \neq m} \frac{1}{m^2 - n^2} = -\frac{3}{4m^2}\) (express the rational fraction \(\frac{1}{m^2 - x^2}\) as a sum of partial fractions).
\item
  If the series whose general term is \(x_{n}\) is convergent in \(\mathbf{R}\) , show that \(\liminf_{n\to \infty}nx_n\leqslant 0\leqslant \limsup_{n\to \infty}nx_n\) .
\end{enumerate}

¶ 10) A series \((u_n)\) is convergent in \(\mathbb{R}\) if and only if, for each increasing sequence \((p_n)\) of numbers \(>0\) such that \(\lim_{n\to \infty}p_n = +\infty\) , we have \(\lim_{n\to \infty}\left(\left(\sum_{k=0}^{n}p_ku_k\right) / p_n\right) = 0\) (use Exercise 15 of § 5).

\begin{enumerate}
\def\labelenumi{\arabic{enumi})}
\setcounter{enumi}{10}
\tightlist
\item
  Consider a sequence \((a_{n})\) every term of which is the sum of a finite sequence of finite real numbers
\end{enumerate}

\[
a _ {n} = b _ {n, 1} + b _ {n, 2} + \dots + b _ {n, k _ {n}}.
\]

For each pair \((n, p)\) of positive integers such that \(p \leqslant k_n\) , if \(m = \sum_{i=0}^{\infty} k_i + p\) , put \(c_m = b_{n,p}\) . Show that if the series whose general term is \(a_n\) is convergent in \(\mathbf{R}\) , and if

\[
d _ {n} = \left| b _ {n, 1} \right| + \left| b _ {n, 2} \right| + \dots + \left| b _ {n, k _ {n}} \right|
\]

tends to o as n tends to infinity, then the series whose general term is \(c_{m}\) is convergent and has the same sum as the series whose general term is \(a_{n}\) .

¶ 12) Let \((u_n)\) be a sequence of finite real numbers. For each permutation \(\sigma\) of the set N, put

\[
r (n) = | \sigma (n) - n |. \sup _ {m \geqslant n} | u _ {m} |.
\]

\begin{enumerate}
\def\labelenumi{\alph{enumi})}
\item
  If the series with general term \(u_n\) is convergent, and if \(\lim_{n \to \infty} r(n) = 0\) , show that the series with general term \(v_n = u_{\sigma(n)}\) is convergent to the same sum. {[}Consider the difference \(\sum_{k=0}^{n} v_k - \sum_{k=0}^{n} u_k\) for large values of \(n\) , and if \(h\) is the smallest integer \(\leqslant n\) such that \(\sigma(h) > n\) , note that in this difference there are at most \(\sigma(h) - h\) terms in each of the two sums which do not cancel out{]}.
\item
  Suppose that the series whose general term is \(u_n\) is convergent. Give an example of a permutation \(\sigma\) such that \(\lim_{n \to \infty} r(n) = 0\) but such that the criterion of Chapter III, § 5, Exercise 6 a) is not satisfied. {[}Define a family of mutually disjoint intervals \(I_k = [n_k - 2k, n_k + 2k]\) in \(N\) , and take \(\sigma\) such that \(\sigma(n) = n\) whenever \(n\) does not belong to any of the \(I_k\) , and such that \(\sigma(n_k - 2j) = n_k + 2j\) , \(\sigma(n_k + 2j) = n_k - 2j\) for all \(k\) and \(0 \leqslant j \leqslant k\) .{]}
\end{enumerate}

\begin{enumerate}
\def\labelenumi{\arabic{enumi})}
\setcounter{enumi}{12}
\item
  Let \((u_n)\) be a sequence of real numbers \(\geqslant 0\) such that \(\lim_{n\to \infty}u_n = 0\) and \(\sum_{n = 0}^{\infty}u_n = +\infty\) . Let \((p_n)\) be a strictly increasing sequence of numbers \(>0\) such that \(\lim_{n\to \infty}p_n = +\infty\) . Show that there exists a permutation \(\sigma\) of \(\mathbf{N}\) such that for each \(n\in \mathbf{N}\) , \(\sum_{k = 0}^{n}u_{\sigma (n)}\leqslant p_n\)
\item
  Let \((u_{n})\) be a sequence of finite real numbers such that the series whose general term is \(u_{n}\) is convergent but not absolutely convergent; let \(s = \sum_{n=0}^{\infty} u_{n}\) . Show that, for each number \(s' \geqslant s\) , there exists a permutation \(\sigma\) of \(\mathbf{N}\) such that \(\sigma(n) = n\) for those indices \(n\) such that \(u_{n} \geqslant 0\) , and such that \(\sum_{n=0}^{\infty} u_{\sigma(n)} = s'\) . {[}Show by induction on \(m\) that there is a permutation \(\sigma_m\) of \(\mathbf{N}\) such that \(\sigma_m(k) = k\) for all \(k\) for which \(u_k \geqslant 0\) and such that, if \(u_k^{(m)}\) denotes \(u_{\sigma_m(k)}\) , there is an integer \(p_m\) with the property that \(\left| s' - \sum_{i=0}^{k} u_i^{(m)} \right| \leqslant 1/m\) for all \(k \geqslant p_m\) . Moreover, \(\sigma_{m+1}\) is such that \(\sigma_{m+1}(k) = \sigma_m(k)\) for all \(k\) such that \(\sigma_m(k) < p_m\) and all \(k\) such that \(u_k \leqslant -1/m\) .{]}
\item
  Let \((u_{n})\) be a sequence of finite real numbers such that \(\lim_{n\to \infty}u_n = 0\) and \(\sum_{n}u_{n}^{+} = \sum_{n}u_{n}^{-} = +\infty\) . If \(a\) and \(b\) are any two real numbers (finite or not) such that \(a\leqslant b\) , show that there is a permutation \(\sigma\) of \(\mathbf{N}\) such that, if we put \(s_n = \sum_{k = 0}^n u_{\sigma (k)}\) , we have \(\liminf_{n\to \infty}s_n = a\) and \(\limsup_{n\to \infty}s_n = b\) . Show that the set of cluster points of the sequence \((s_n)\) is then the interval \([a,b]\) .
\item
  Let \((u_n)\) be a sequence of finite real numbers such that the series whose general term is \(u_n\) is convergent but not absolutely convergent. Show that there exists a permutation \(\sigma\) of \(\mathbf{N}\) satisfying the condition of Chapter III, § 5, Exercise 6 \(a\) ) but such that the series whose general term is \(u_{\sigma - 1(n)}\) is not convergent. {[}Let \(h, k, m\) be three integers with the following properties: if \(p_1, \ldots, p_r\) are the integers \(n\) such that \(h \leqslant n < h + m\) and \(u_n \geqslant 0\) , then \(u_{p_1} + \cdots + u_{p_r} > 2\) and \(h + m < k - r\) ; moreover if \(\mu, \nu\) are any integers such that \(h + m \leqslant \mu \leqslant \nu\) , then \(\left|\sum_{n = \mu}^{\nu} u_n\right| \leqslant 1\) . Put \(s = m - r\) and let \(q_1, \ldots, q_s\) denote the integers \(n\) such that \(h \leqslant n < h + m\) and \(u_n < 0\) . Consider the permutation \(\pi\) of the interval \([h, k + s]\) of \(\mathbf{N}\) such that \(\pi(p_i) = k - i + 1\) for \(i \leqslant i \leqslant r\) , \(\pi(q_j) = k + j\) for \(i \leqslant j \leqslant s\) , \(\pi(k + j) = h + s - j\) for \(i \leqslant j \leqslant s\) , \(\pi(k - i + 1) = h + s + i - 1\) for \(i \leqslant i \leqslant r\) and finally \(\pi(n) = n\) for \(h + m \leqslant n \leqslant k - r\) ; show that we have
\end{enumerate}

\[
\left| \sum_ {i = h} ^ {k} u _ {i} - \sum_ {i = h} ^ {k} u _ {\pi^ {- 1} (i)} \right| \geqslant I.
\]

\begin{enumerate}
\def\labelenumi{\arabic{enumi})}
\setcounter{enumi}{16}
\tightlist
\item
  If \(u_{mn} = \mathbf{i} / (m^2 - n^2)\) for \(m \neq n\) and \(u_{mm} = 0\) , show that
\end{enumerate}

\[
\sum_ {m = 1} ^ {\infty} \left(\sum_ {n = 1} ^ {\infty} u _ {m n}\right) = - \sum_ {n = 1} ^ {\infty} \left(\sum_ {m = 1} ^ {\infty} u _ {m n}\right) \neq 0.
\]

(Use Exercise 8.)

\begin{enumerate}
\def\labelenumi{\arabic{enumi})}
\setcounter{enumi}{17}
\item
  Let \((\mathbf{i} + u_{\mathfrak{i}})\) be a family of real numbers all belonging to the interval \([0, +\infty[\) (resp. \(]0, +\infty]\) ). Then the family \((\mathbf{i} + u_{\mathfrak{i}})\) is multipliable in \(\mathbb{R}\) if and only if one of the following conditions is satisfied: a) At least one of the numbers \(\mathbf{i} + u_{\mathfrak{i}}\) is equal to o (resp. \(+\infty\) ). b) The family \(u_{\mathfrak{i}}\) is summable in \(\overline{\mathbb{R}}\) .
\item
  Let \((\mathbf{I} + u_{t})_{t\in \mathbf{I}}\) be a multipliable family in \(\mathbb{R}\) . For each non-empty subset \(\mathbf{H}\) of \(\mathbf{I}\) , let \(v_{\mathbf{H}}\) denote \(\prod_{i\in \mathbf{H}}u_i\) . Show that the family \((v_{\mathbf{H}})_{\mathbf{H}\in \mathfrak{F}(\mathbf{I})}\) is summable in \(\mathbb{R}\) and that
\end{enumerate}

\[
\mathrm{I} + \sum_ {\mathbf {H}} v _ {\mathbf {H}} = \prod_ {i \in \mathbf {I}} (\mathrm{I} + u _ {i}).
\]

Is the converse true?

Deduce that, if \(-\mathbf{I} < x < \mathbf{I}\) , then \(\prod_{n=0}^{\infty} (\mathbf{I} + x^{2^n}) = \mathbf{I} / (\mathbf{I} - x)\) .

\begin{enumerate}
\def\labelenumi{\arabic{enumi})}
\setcounter{enumi}{19}
\item
  Let \((u_n)\) be a sequence of finite real numbers \(\geqslant 0\) such that \(u_0 > 0\) , and let \(s_n = \sum_{k=0}^{n} u_k\) for each \(n \geqslant 0\) . Then the sequence \((u_n)\) is summable in \(\mathbf{R}\) if and only if the sequence \((u_n / s_n)\) is summable in \(\mathbf{R}\) {[}apply Theorem 4 to the sequence \((1 - u_n / s_n)]\) . The same holds if the sequence \((u_n / s_n)\) is replaced by the sequence \((u_n / s_{n-1})\) .
\item
  Let \(u_{n} = (-\mathrm{I})^{n} / \sqrt{n}\) for \(n \geqslant 2\) . The product whose general factor is \(\mathrm{I} + u_{n}\) is not convergent, but the series whose general term is \(u_{n}\) is convergent.
\item
  For \(n \geqslant 2\) , let \(u_{2n-1} = -\mathrm{I}/\sqrt{n}\) and \(u_{2n} = (\mathrm{I} + \sqrt{n})/n\) . The product whose general factor is \(\mathrm{I} + u_n\) is convergent, but the series whose general term is \(u_n\) is not convergent.
\end{enumerate}

\BourbakiExerciseGroup

\begin{enumerate}
\def\labelenumi{\arabic{enumi})}
\tightlist
\item
  If \(x\) and \(y\) are real, we have
\end{enumerate}

\[
[ x + y ] = [ x ] + [ y ] + \varepsilon \quad \text { where } \varepsilon = 0 \quad \text { or } \quad \varepsilon = 1,
\]

\[
[ x - y ] = [ x ] - [ y ] - \varepsilon \quad \text { where } \varepsilon = 0 \quad \text { or } \quad \varepsilon = 1,
\]

\[
[ x ] + [ x + y ] + [ y ] \leqslant [ 2 x ] + [ 2 y ],
\]

\[
[ x ] + \left[ x + \frac {1}{n} \right] + \left[ x + \frac {2}{n} \right] + \dots + \left[ x + \frac {n - 1}{n} \right] = [ n x ],
\]

\[
\left[ \frac {[ n x ]}{n} \right] = [ x ]
\]

for any integer \(n > 0\) .

\begin{enumerate}
\def\labelenumi{\arabic{enumi})}
\setcounter{enumi}{1}
\item
  Let \((\varepsilon_{n})\) be a strictly decreasing sequence of finite numbers \(>0\) which tend to 0. For each \(x \in \mathbb{R}\) , let \(r_{n}(x)\) be the multiple of \(\varepsilon_{n}\) which is the approximation by defect to \(x\) to within \(\varepsilon_{n}\) . Then the sequence \((r_{n}(x))\) is increasing for all \(x \in \mathbb{R}\) if and only if \(\varepsilon_{n}\) is an integer multiple of \(\varepsilon_{n+1}\) for each \(n\) .
\item
  Let \(a\) be an integer \(> 1\) . A real number \(x\) is rational if and only if its expansion to base \(a\) , say \(x = p_0 + \sum_{n=0}^{\infty} u_n a^{-n}\) , is periodic, that is to say if and only if there exist two integers \(n_0\) and \(r > 0\) such that \(u_{n+r} = u_n\) for all \(n \geqslant n_0\) .
\item
  Let \((u_{n})\) be a summable sequence of numbers \(>0\) in \(\mathbf{R}\) which satisfies the following conditions: \(u_{n+1} \leqslant u_{n}\) and \(u_{n} \leqslant \sum_{k=1}^{\infty} u_{n+k}\) for all \(n \geqslant 0\) . Show that for every number \(a\) such that \(0 < a \leqslant s = \sum_{n=0}^{\infty} u_{n}\) , there exists a subset I of \(\mathbf{N}\) such that \(a = \sum_{n \in \mathbf{I}} u_{n}\) . These conditions are satisfied in particular if \(u_{n+1} \leqslant u_{n} \leqslant 2u_{n+1}\) for all \(n\) . Consider the case of dyadic expansions.
\item
  For each real number \(x \in ]0, I]\) , there is a unique infinite increasing sequence \((q_n)\) of integers \(> 0\) such that \(x = \sum_{n=1}^{\infty} I / (q_1 q_2 \ldots q_n)\) . The number \(x\) is rational if and only if \(q_{n+1} = q_n\) from a certain value of \(n\) onwards.
\item
  For each real number \(x > \mathbf{I}\) , there is a unique infinite sequence \((t_n)\) of integers \(\geqslant \mathbf{I}\) such that \(t_{n+1} \geqslant t_n^2\) for all \(n \geqslant \mathbf{I}\) and such that \(x = \prod_{n=1}^{\infty} \left( \mathbf{I} + \frac{\mathbf{I}}{t_n} \right)\) (use Exercise 19 of §7). The number \(x\) is rational if and only if \(t_{n+1} = t_n^2\) form a certain value of \(n\) onwards. {[}To show that the condition is necessary, show that if we put \(x_n = \prod_{k=n+1}^{\infty} \left( \mathbf{I} + \frac{\mathbf{I}}{t_k} \right)\) then the denominators of the fractions \(\frac{x_n}{x_n - \mathbf{I}}\) (in their irreducible form) form a decreasing sequence{]}.
\end{enumerate}

¶ * 7) Let \((\varepsilon_n)_{n \geqslant 0}\) be a sequence of numbers each of which is 1 or --- 1. Show that the real number

\[
x _ {n} = \varepsilon_ {0} \sqrt {2 + \varepsilon_ {1} \sqrt {2 + \varepsilon_ {2} \sqrt {2 + \cdots + \varepsilon_ {n} \sqrt {2}}}}
\]

is defined and is equal to

\[
2 \sin \left(\frac {\pi}{4} \sum_ {k = 0} ^ {n} \frac {\varepsilon_ {0} \varepsilon_ {1} \dots \varepsilon_ {k}}{2 ^ {k}}\right).
\]

Deduce that \(\lim_{n\to \infty}x_n\) exists and that for each real number \(x\) , such that \(-2\leqslant x\leqslant 2\) , there is a sequence \((\varepsilon_n)\) such that \(x\) is equal to the limit of the corresponding sequence \((x_{n})\) . For what values of \(x\) is the sequence \((\varepsilon_n)\) unique? For what values of \(x\) is it periodic? (see Exercise 3).*

\begin{enumerate}
\def\labelenumi{\arabic{enumi})}
\setcounter{enumi}{7}
\item
  Show that there is no non-constant mapping \(\psi\) of an interval \(\mathbf{I} \subset \mathbf{R}\) into the set \(\mathbf{N}^{\mathbb{N}}\) of all sequences of natural numbers which has the following property for each \(x \in \mathbf{I}\) : for each integer \(n \geqslant 0\) there is a neighbourhood \(\mathbf{V}\) of \(x\) in \(\mathbf{I}\) such that, for each \(y \in \mathbf{V}\) , the first \(n\) terms of the sequence \(\psi(y)\) are the same as the first \(n\) terms of the sequence \(\psi(x)\) (remark that this would imply the continuity of \(\psi\) , if each factor \(\mathbf{N}\) carries the discrete topology and \(\mathbf{N}^{\mathbb{N}}\) the product topology).
\item
  Show that Cantor's triadic set K (§ 2, no. 5) is the set of all real numbers \(x \in [0, 1]\) whose triadic expansion (resp. improper triadic expansion if \(x\) is the origin of an interval contiguous to K)
\end{enumerate}

\[
x = \sum_ {n = 1} ^ {\infty} u _ {n} 3 ^ {- n}
\]

is such that \(u_{n} \neq \mathrm{I}\) for each \(n\) (so that \(u_{n} = 0\) or \(u_{n} = 2\) ). If \(v_{n} = \frac{\mathrm{I}}{2} u_{n}\) for each \(n\) , and if \(f(x) = \sum_{n=1}^{\infty} v_{n} 2^{-n}\) , show that \(f\) is a surjective continuous mapping of \(\mathbf{K}\) onto \([\mathrm{o}, \mathrm{I}]\) ; deduce that \(\mathbf{K}\) has the power of the continuum.

¶ 10) Let \((d_n)\) be a base sequence, let \(a_n = d_n / d_{n-1}\) ( \(n \geqslant 1\) ), let \((n_i)\) be a strictly increasing sequence of integers, and let \((b_i)\) be a sequence of integers such that \(0 < b_i < a_{n_i} - 1\) for each \(i\) . Let \(G\) be the set of all \(x \in [0, 1]\) such that, if \(\sum_{n=1}^{\infty} u_n / d_n\) is the expansion of \(x\) , or the improper expansion if it exists, then \(u_{n_i} \neq b_i\) for all \(i \in N\) . Show that \(G\) is a totally disconnected perfect set, and that if \(l\) is the sum of the lengths of the intervals contiguous to \(G\) and contained in \([0, 1]\) , then

\[
\mathrm{I} - l = \prod_ {i = 0} ^ {\infty} \left(\mathrm{I} - \frac {\mathrm{I}}{a _ {n _ {i}}}\right).
\]

¶ 11) Let X be a Hausdorff topological space. Suppose that, for each finite sequence s whose terms are equal to o or 1, there is a nonempty subset \(\mathrm{A}(s)\) of X which satisfies the following conditions: (i) If s is a sequence of n terms (each of which is o or i) and \(s'\) , \(s''\) are the sequences of \(n + 1\) terms (each of which is o or i) whose first n terms are the same as those of s, then \(\mathrm{A}(s) = \mathrm{A}(s') \cup \mathrm{A}(s'')\) ; and if \(s_0\) is the empty sequence, \(\mathrm{A}(s_0) = \mathrm{X}\) .

\begin{enumerate}
\def\labelenumi{(\roman{enumi})}
\setcounter{enumi}{1}
\tightlist
\item
  For each infinite sequence \((u_n)_{n\geqslant 0}\) whose terms are equal to 0 or 1, if \(s_n\) denotes the finite sequence \((u_k)_{0\leqslant k\leqslant n}\) , then the filter base formed by the \(\mathbf{A}(s_n)\) converges to a point of \(\mathbf{X}\) .
\end{enumerate}

Show that under these conditions there is a surjective continuous mapping of Cantor's triadic set K onto X.

¶ 12) Deduce from Exercise 11 that:

\begin{enumerate}
\def\labelenumi{\alph{enumi})}
\item
  If A is a compact subset of R, there is a continuous mapping of Cantor's triadic set K onto A {[}take the A(s) to be the intersections of A with suitably chosen intervals{]}.
\item
  If also A is perfect and totally disconnected, it is homeomorphic to K {[}same method, by arranging that if \(s_{1}\) and \(s_{2}\) are two distinct sequences with the same number of terms (equal to o or 1) then
\end{enumerate}

\[
\mathbf {A} (s _ {1}) \cap \mathbf {A} (s _ {2}) = \varnothing_ {1} ].
\]

¶ 13) Let X be a countable Hausdorff space. Suppose that, for each finite sequence s all of whose terms are o or 1, there is a non-empty subset B(s) of X which satisfies the following conditions:

\begin{enumerate}
\def\labelenumi{(\roman{enumi})}
\item
  If \(s\) is a finite sequence of \(n\) terms (equal to \(\mathbf{o}\) or \(\mathbf{i}\) ) and \(s', s''\) are the sequences of \(n + \mathbf{i}\) terms (equal to \(\mathbf{o}\) or \(\mathbf{i}\) ) whose first \(n\) terms are the same as those of \(s\) , then \(\mathbf{B}(s') \cup \mathbf{B}(s'') = \mathbf{B}(s)\) and \(\mathbf{B}(s') \cap \mathbf{B}(s'') = \emptyset\) ; and if \(s_0\) is the empty sequence, \(\mathbf{B}(s_0) = \mathbf{X}\) .
\item
  For each \(x \in \mathbf{X}\) , if \(s_n\) is the (unique) sequence of \(n\) terms equal to 0 or 1 such that \(x \in \mathrm{B}(s_n)\) , then the filter base formed by the \(\mathrm{B}(s_n)\) converges to \(x\) in \(\mathbf{X}\) .
\end{enumerate}

Show that, under these conditions, E is homeomorphic to the rational line Q. {[}Show first that X is homeomorphic to a countable subset of Cantor's triadic set K which is dense in K and contains no end-point of an interval contiguous to K. In order to do this, remark that the hypotheses associate with each \(x \in X\) an infinite sequence \((u_{n}(x))\) whose terms are equal to 0 or 1; using the fact that X is countable, show that the bijection \(s \to \mathbf{B}(s)\) can be modified in such a way that for no \(x \in E\) do the \(u_{n}(x)\) form a stationary sequence; finally use Exercise 10 of §2{]}.

Deduce that every countable subspace of \(\mathbf{R}\) which has no isolated points is homeomorphic to \(\mathbf{Q}\) .

\begin{enumerate}
\def\labelenumi{\arabic{enumi})}
\setcounter{enumi}{13}
\tightlist
\item
  Show that every closed subset of \(\mathbf{R}\) is either countable or has the power of the continuum {[}use Exercise 12 b) above and Exercise 16 of Chapter I, § 9{]}.
\end{enumerate}


\begin{enumerate}
\def\labelenumi{\arabic{enumi})}
\setcounter{enumi}{14}
\item
  Show that the set of open subsets of \(\mathbf{R}\) and the set of compact, totally disconnected, perfect subsets of \(\mathbf{R}\) have the power of the continuum.
\item
  \begin{enumerate}
  \def\labelenumii{\alph{enumii})}
  \tightlist
  \item
    Let A be a countable closed subset of R and let f be a continuous real-valued function defined on R. Show that if f is constant in each of the intervals contiguous to A, then f is constant (use Bolzano's theorem).
  \end{enumerate}
\end{enumerate}

\begin{enumerate}
\def\labelenumi{\alph{enumi})}
\setcounter{enumi}{1}
\tightlist
\item
  If \(f\) is the continuous mapping of Cantor's triadic set \(\mathbf{K}\) onto \([0, 1]\) defined in Exercise 9, show that \(f\) can be extended to a continuous function on \(\mathbf{R}\) which is constant on each of the intervals contiguous to \(\mathbf{K}\) .
\end{enumerate}

\begin{enumerate}
\def\labelenumi{\arabic{enumi})}
\setcounter{enumi}{16}
\tightlist
\item
  Let X be a topological space which has a countable dense subset. Show that the set of all real-valued continuous functions on X has the power of the continuum.
\end{enumerate}
